\section{Signed curvature separation and small loops}
\label{holo:section}

Fix a datum $d$, a noncorner branch seed $w$, and a finite data
order $m\geq10$.  We use the signed experiments of
Proposition~\ref{energy:signed-experiment} and their attachment and
observers in Propositions~\ref{attach:development}
and~\ref{attach:observer}.  The observation point is $p_T$, with
$v(p_T)=T$ and
\[
 a_p=a(p_T)\asymp e^{-\kappa T},\qquad
 A_T=e^{-7\kappa T/8},\qquad \lambda_T=e^{\zeta T}.
\]
All estimates below concern a coordinate box of fixed positive size
about $p_T$, contained with strict margins in the attached future
development.  On this box $a/a_p$ and $a_p/a$ are bounded.  A member
of a fixed finite angular atlas contains a smaller such box; its
constants can be maximized over the atlas.  All derivatives denoted
by $D$ are ordinary coordinate derivatives in this box.

Let $k_T^\pm=g_T^\pm-g$ be the physical metric differences, and let
$h_T^{\mathrm p}$ be the physical linear packet.  Their arrays in
the background normalized frame are respectively $H_T^\pm$ and
$q_T$.  Thus the arrays of $k_T^\pm\mp h_T^{\mathrm p}$ are
$F_T^\pm$, and $F_T^{\mathrm{odd}}=F_T^+-F_T^-$.  This distinguishes
a tensor from its relative array whenever slots are evaluated.
The estimates already proved imply, for some $c>0$,
\begin{align}
 \|H_T^\pm\|_{C^{30}}&\leq
       e^{-(\kappa/2+c)T},                                      \label{holo:relative-small}\\
 e^{o(T)}\bigl(\|F_T^{\mathrm{odd}}\|_{C^2}
       +\|H_T^+\|_{C^2}^2+\|H_T^-\|_{C^2}^2\bigr)&=o(A_T).
                                                               \label{holo:odd-small}
\end{align}
Here and below $e^{o(T)}$ refers to the fixed finite orders in use.
These two estimates have different roles. The bound $H_T^\pm=o(\sqrt a)$
implied by Equation~\eqref{holo:relative-small} controls the connection
after passing to the observer frame, just as it controlled the terminal
state conversion in Section~\ref{attach:section}. Equation~\eqref{holo:odd-small}
makes the odd error and the quadratic curvature remainder smaller than
the linear amplitude $A_T$. It follows from
Equation~\eqref{energy:small-relative}, including its fixed-order
subexponential factors. The individual errors $F_T^\pm$ need not be
smaller than the packet.

Choose a smooth $\gamma$-orthonormal angular frame $e_2,e_3$, with
$e_2(p_T)=X_T$, the polarization direction in
Proposition~\ref{bg:packet}.  Gram--Schmidt in the fixed atlas and a
constant orthogonal rotation give ordinary jets of size $e^{o(T)}$.
Put $\hat e_i=a^{-1/2}L_i$ for $i=0,1$ and $\hat e_A=e_A$ for
$A=2,3$.  Its Lorentz Gram matrix $\eta$ is constant:
$\eta_{01}=\eta_{10}=-1$, $\eta_{22}=\eta_{33}=1$, and the other
entries vanish.  Re-express $H_T^\pm,q_T,F_T^\pm$ in this angular
frame, retaining their notation.  The frame change and its inverse
have subexponential fixed jets, so
Equations~\eqref{holo:relative-small}--\eqref{holo:odd-small}
remain valid, with $c$ decreased if necessary.  Define
\begin{equation}\label{holo:frames}
 I_T^\pm=(g^{-1}g_T^\pm)^{-1/2},\qquad
 E_0^\pm=I_T^\pm L_0,\quad
 E_1^\pm=a^{-1}I_T^\pm L_1,\quad
 E_A^\pm=I_T^\pm e_A.
\end{equation}
The inverse square root is the analytic branch at the identity.
It is an isometry from $g$ to $g_T^\pm$, so these frames also have
Gram matrix $\eta$.  By Proposition~\ref{attach:observer}, their
Euclidean component norms at $p_T$ are uniformly comparable to the
positive test norm for the adaptive seeds $w_T^\pm\to w$.

\subsection{Separation of a scaled curvature component}

We first take the difference of the two scaled curvature observations,
then read the selected observation from finite jets of a family of loops.
The geometric input is Proposition~\ref{bg:packet}, whose proof requires
no bound on the observer's background curvature.

\begin{lemma}[Smooth scaled curvature dependence]
\label{holo:curvature-functional}
At $p_T$ define
\begin{equation}\label{holo:scaled-curvature}
 C_T^\pm=a_p^2\Riem(g_T^\pm)
       (E_1^\pm,E_2^\pm,E_1^\pm,E_2^\pm).
\end{equation}
There is a smooth function $\mathcal C_T$ of the relative metric
array and its ordinary jets through order two such that
$C_T^\pm=\mathcal C_T(H_T^\pm)$.  On a fixed small neighborhood
of the zero relative array its first two derivatives have norm
at most $e^{o(T)}$.  Consequently
\begin{equation}\label{holo:signed-taylor}
 C_T^+-C_T^-=2D\mathcal C_T(0)[q_T]+o(A_T).
\end{equation}
\end{lemma}

\begin{proof}
We give the scaling calculation behind the assertion.  Set
$Z_\alpha=\sqrt a\,\hat e_\alpha$, and let $V_\alpha$ be the
matrices of $\nabla^g_{Z_\alpha}$ on the $\hat e$ frame.  The
coefficient estimates of Lemma~\ref{bg:coefficients} bound $V$,
the coefficients of $[Z_\alpha,Z_\beta]$ in $Z$, their required
ordinary jets, and the normalized entries of $a\Riem(g)$ by
$e^{o(T)}$.  Relative derivatives of $a$ have the same bound.
These are the coefficient inputs in \cite[\CompanionRef{foundation}{pk2:coefficients}]{Foundation}.

For an arbitrary small relative change $H$, write $M=\eta+H$ and
let $k$ be the physical tensor with relative array $H$.
The connection difference $K=\nabla^{g+k}-\nabla^g$ satisfies
\[
 2(g+k)(K(V,W),Y)
    =(\nabla^g_Vk)(W,Y)+(\nabla^g_Wk)(V,Y)
                              -(\nabla^g_Yk)(V,W).
\]
On normalized slots, multiplication by $\sqrt a$ replaces each
derivative on the right by a $Z$ derivative.  More explicitly, put
\[
 \mathcal D_\alpha H_{\beta\delta}
   =Z_\alpha H_{\beta\delta}
      -(V_\alpha)^\rho{}_{\beta}H_{\rho\delta}
      -(V_\alpha)^\rho{}_{\delta}H_{\beta\rho}.
\]
The array $W$ of $\sqrt a K$ is therefore
\begin{equation}\label{holo:connection-difference}
 W^\mu{}_{\alpha\beta}
   =\frac12(M^{-1})^{\mu\delta}
      \bigl(\mathcal D_\alpha H_{\beta\delta}
          +\mathcal D_\beta H_{\alpha\delta}
          -\mathcal D_\delta H_{\alpha\beta}\bigr).
\end{equation}
Thus $V+W$ is the connection of $g+k$, differentiated along $Z$
and expressed in the unchanged normalized frame.  If
$[Z_\alpha,Z_\beta]=b^\rho{}_{\alpha\beta}Z_\rho$, its curvature
matrix is
\begin{equation}\label{holo:scaled-curvature-formula}
 R^{g+k}(Z_\alpha,Z_\beta)
  =Z_\alpha(V_\beta+W_\beta)-Z_\beta(V_\alpha+W_\alpha)
   +[V_\alpha+W_\alpha,V_\beta+W_\beta]
   -b^\rho{}_{\alpha\beta}(V_\rho+W_\rho).
\end{equation}
Its entries are smooth functions of $H,DH,D^2H$ with
subexponential coefficients, since $Z_i=L_i$ and $Z_A=\sqrt a e_A$
have such coordinate jets.  Contracting with $M$ gives every
normalized covariant entry of $a\Riem(g+k)$.  The isometry $I(H)$
is likewise an analytic function of $\eta^{-1}H$; hence including
its four slots preserves these bounds.

At $p_T$, Equation~\eqref{holo:scaled-curvature} equals $a_p$
times the curvature on slots
$I(H)\hat e_1,I(H)e_2,I(H)\hat e_1,I(H)e_2$.
Equations~\eqref{holo:connection-difference} and
\eqref{holo:scaled-curvature-formula} therefore prove the claimed
jet dependence and derivative bounds.  This controls the scaled
background coefficients; it does not assert bounded curvature in
the boosted observer frame.

Taylor's formula at zero, applied separately to the two signs,
gives
\[
 C_T^+-C_T^-
 =2D\mathcal C_T(0)[q_T]
       +D\mathcal C_T(0)[F_T^{\mathrm{odd}}]+\mathcal E_T,
 \qquad
 |\mathcal E_T|\leq e^{o(T)}
       (\|H_T^+\|_{C^2}^2+\|H_T^-\|_{C^2}^2).
\]
Equation~\eqref{holo:odd-small} proves
Equation~\eqref{holo:signed-taylor}.  The common term
$\mathcal C_T(0)$ has canceled exactly.
\end{proof}

\begin{proposition}[Signed signal]
\label{holo:signal}
For every sufficiently large $T$, at least one sign satisfies
$|C_T^\pm|>A_T$.
\end{proposition}

\begin{proof}
We compute the linear term in Equation~\eqref{holo:signed-taylor}.
Up to the curvature sign convention, the principal second derivative
part of the variation on the physical slots $L_1,X_T,L_1,X_T$ is
\[
 \frac12\bigl(\nabla_1\nabla_1 h_{XX}
              +\nabla_X\nabla_X h_{11}
              -2\nabla_1\nabla_X h_{1X}\bigr).
\]
The packet from Proposition~\ref{bg:packet} is
\[
 h_T^{\mathrm p}
   =A_T\operatorname{tr.rev.}_g\operatorname{Re}
       \left(e^{i\lambda_T v}
                    \sum_{j=0}^{N}\lambda_T^{-j}p_j\right).
\]
The finite $N$ has already been fixed.  Each amplitude has fixed
ordinary normalized jets of size $e^{o(T)}$; $p_0$ is trace-free
and annihilates $L_0$.  Since $dv(L_1)=1$ and $dv(X_T)=0$, only
the two phase derivatives on the $XX$ entry contribute order
$A_T\lambda_T^2$.  The leading trace reversal has no effect.
All terms with a derivative on an amplitude or connection, all
$j\geq1$ terms, and the variation of the four isometry slots on
the scaled background curvature have at most order
$A_T\lambda_T e^{o(T)}$.  The last assertion also follows directly
from Equation~\eqref{holo:scaled-curvature-formula}: differentiation
of $I(H)$ is algebraic and so creates no phase derivative.
We obtain
\begin{equation}\label{holo:linear-signal}
 D\mathcal C_T(0)[q_T]
   =\pm\frac12 A_T\lambda_T^2
       \operatorname{Re}\bigl(e^{i\lambda_TT}p_0(X_T,X_T)\bigr)
       +O(A_T\lambda_T e^{o(T)}).
\end{equation}
The choice of direction and constant complex phase in
Proposition~\ref{bg:packet} makes the real contraction here equal to a
positive $\sigma_T$ with $\sigma_T$ and $\sigma_T^{-1}$ of size
$e^{o(T)}$.  The source construction's observer-orthogonal screen
representative differs from the angular $X_T$ by a multiple of
$L_0$; its contraction with $p_0$ is identical because
$p_0(L_0,\cdot)=0$.  Thus the polarization input precedes, and does
not use, the bounded-curvature hypothesis in the tidal application
of~\cite[\CompanionRef{foundation}{pk2:deep-packet}]{Foundation}.
Since $\lambda_T=e^{\zeta T}$ with fixed $\zeta>0$,
\[
 \frac{A_T\lambda_T e^{o(T)}}{A_T\lambda_T^2\sigma_T}
      =e^{-\zeta T+o(T)}\longrightarrow0,
 \qquad \lambda_T^2\sigma_T\longrightarrow\infty.
\]
Equation~\eqref{holo:signed-taylor} now gives
$|C_T^+-C_T^-|>2A_T$ for large $T$, proving the claim by the
triangle inequality.
\end{proof}

\subsection{Connection bounds in the observer frame}

The loop estimates will be used to recover a quadratic holonomy
coefficient from finitely many small values. If its eighteenth
parameter derivative is controlled, interpolation gives an upper
bound of size $e^{o(T)}(a_p/h+h^{16})$. We will take
$h=e^{-\kappa T/16}$: both terms are then $o(A_T)$, since
$a_p\asymp e^{-\kappa T}$ and $A_T=e^{-14\kappa T/16}$.
This is the reason for controlling eighteen holonomy derivatives;
the twenty connection derivatives below leave the needed ODE buffer.
We first show that the large boost in the observer frame does not
spoil those bounds.

Fix a sign given by Proposition~\ref{holo:signal}, and abbreviate
its metric, isometry, and frame by $g'$, $I$, and $E$.  Set
\begin{equation}\label{holo:loop-frame}
 f_\alpha=a_pE_\alpha,\qquad
 \nabla^{g'}_{f_\alpha}E_\beta
        =G^\mu{}_{\alpha\beta}E_\mu.
\end{equation}
The number $a_p$ is held constant on the box.

\begin{lemma}[Connection and bracket jets]
\label{holo:boosted-control}
The coordinate components of $f_\alpha$, the arrays
$G^\mu{}_{\alpha\beta}$, and the coefficients
$B^\mu{}_{\alpha\beta}$ of $[f_\alpha,f_\beta]$ in the $f$
frame have ordinary jets through order twenty bounded by $e^{o(T)}$.
\end{lemma}

\begin{proof}
Write $\widetilde e_\alpha=I\hat e_\alpha$ and let
$\widetilde V^\mu{}_{\alpha\beta}$ be defined by
\[
 \nabla^{g'}_{\sqrt a\,\widetilde e_\alpha}
                   \widetilde e_\beta
     =\widetilde V^\mu{}_{\alpha\beta}\widetilde e_\mu.
\]
Equation~\eqref{holo:connection-difference}, the change of frame
formula, and Equation~\eqref{holo:relative-small} give
\begin{equation}\label{holo:perturbed-connection-small}
 \|\widetilde V-V\|_{C^{20}}
       +\|I-\Id\|_{C^{21}}
       \leq e^{-(\kappa/2+c_1)T}
\end{equation}
for some $c_1>0$, with $I-\Id$ expressed in normalized entries.
Indeed the change of frame formula uses $I,I^{-1},Z I$ and the
matrices $V+W$; every difference has a factor $H$ or $DH$.
Twenty derivatives use at most twenty-one metric derivatives;
the thirty-derivative input leaves a buffer.  Multiplication by
any fixed collection of background $e^{o(T)}$ factors only reduces
$c$ to a smaller positive $c_1$.

Put
\[
 s_0=\tfrac12,\qquad s_1=-\tfrac12,\qquad s_2=s_3=0,
 \qquad E_\alpha=a^{s_\alpha}\widetilde e_\alpha.
\]
The full conversion formula is
\begin{equation}\label{holo:boost-conversion}
 G^\mu{}_{\alpha\beta}
   =\frac{a_p}{a}
      a^{1/2+s_\alpha+s_\beta-s_\mu}
                   \widetilde V^\mu{}_{\alpha\beta}
       +s_\beta(f_\alpha\log a)\delta^\mu_\beta.
\end{equation}
There are precisely seven triples with negative exponent in the
first term.  The following table lists them all; $A=2,3$ in each
of its last three rows.
\begin{center}
\begin{tabular}{c c c l}
$(\alpha,\beta,\mu)$ & exponent & number & reason for control\\
\hline
$(1,1,0)$ & $-1$ & $1$ & exact null metric compatibility\\
$(A,1,0)$ & $-1/2$ & $2$ & exact null metric compatibility\\
$(1,1,A)$ & $-1/2$ & $2$ & background null geodesic identity\\
$(1,A,0)$ & $-1/2$ & $2$ & background angular orthogonality
\end{tabular}
\end{center}
For completeness, $s_\alpha+s_\beta$ belongs to
$\{-1,-1/2,0,1/2,1\}$.  If $\mu=1$ the exponent is
$1+s_\alpha+s_\beta\geq0$.  If $\mu$ is angular it is negative
only when $\alpha=\beta=1$.  If $\mu=0$ it equals
$s_\alpha+s_\beta$, negative only for $(1,1),(1,A),(A,1)$.
This proves exhaustiveness.

Nullness gives, for every $\alpha$,
\[
 0=\tfrac12\sqrt a\,\widetilde e_\alpha
                   g'(\widetilde e_1,\widetilde e_1)
   =g'(\nabla^{g'}_{\sqrt a\widetilde e_\alpha}
                     \widetilde e_1,\widetilde e_1)
   =-\widetilde V^0{}_{\alpha1}.
\]
Thus the first three individual entries vanish identically for the
perturbed metric, together with every derivative.  In the other
four entries only the background vanishes.  In fact
$L_1=\partial_v$ in the optical metric satisfies
$\nabla^g_{L_1}L_1=(L_1\log a)L_1$, so
$V^A{}_{11}=0$.  Differentiating
$g(\hat e_1,e_A)=0$ along $L_1$ then gives
$V^0{}_{1A}=0$.  These are identities throughout the box, hence
their ordinary derivatives also vanish.  Equation~\eqref{holo:perturbed-connection-small}
therefore bounds all four remaining perturbed entries.  Their
worst multiplier is $a^{-1/2}$, leaving
\[
 a^{-1/2}e^{-(\kappa/2+c_1)T}=O(e^{-c_1T})
\]
on this box.  Derivatives of the boost factors retain these powers,
because $D^j a^r=a^r e^{o(T)}$ at every fixed order.  Thus all
twenty derivatives of the seven entries are controlled.

We also verify the coordinate bounds used in the diagonal term of
Equation~\eqref{holo:boost-conversion}.  A normalized radial basis
vector has coordinate size $a^{-1/2}e^{o(T)}$, whereas an angular
one has size $e^{o(T)}$.  For a term with input label $\alpha$
and normalized output label $\rho$, the largest coordinate
multiplier is $a_pa^{s_\alpha-1/2}$ if $\rho$ is radial, and
$a_pa^{s_\alpha}$ if it is angular.  As $a\asymp a_p$, their
worst powers are respectively $a^0$ and $a^{1/2}$, both attained
when $\alpha=1$.  Consequently the coordinates of $f$ and their
twenty derivatives are $e^{o(T)}$.  The same holds for
$f_\alpha\log a$ by the relative derivative bounds for $a$.
The nonnegative exponents in Equation~\eqref{holo:boost-conversion}
are now controlled as well.

Finally, torsion-freeness and constancy of $a_p$ give the exact
identity
\begin{equation}\label{holo:bracket}
 [f_\alpha,f_\beta]
    =(G^\mu{}_{\alpha\beta}-G^\mu{}_{\beta\alpha})f_\mu.
\end{equation}
There is no factor $a_p^{-1}$.  This proves the bracket jet bounds
without estimating the inverse coordinate matrix of $f$.
\end{proof}

\subsection{The parameter-dependent flow and its developed length}

Set
\begin{equation}\label{holo:h-choice}
 h_T=e^{-\kappa T/16}.
\end{equation}
For $0\leq h\leq17h_T$, let
\begin{equation}\label{holo:surface}
 \Psi_h(s,t)=\operatorname{Fl}^{1}_{h(sf_1+tf_2)}(p_T),
                  \qquad (s,t)\in[0,1]^2.
\end{equation}
The vector field in this expression depends on the parameters
$s,t,h$, and $\operatorname{Fl}^1$ is its actual time-one flow.
Let $\ell_h$ be the image of the positively oriented square
boundary, based at $(0,0)$, and let $\mathcal H_T(h)$ be its
parallel transport matrix in $E(p_T)$.  Write $\mathcal L_T(h)$
for its length in the positive metric obtained by continuing the
test's parallel metric along the loop from $p_T$.

\begin{lemma}[Loop estimates]
\label{holo:loops}
For sufficiently large $T$, every loop in
Equation~\eqref{holo:surface} is defined inside the attached future
observation box.  There is $Q_T=e^{o(T)}$ such that
\begin{equation}\label{holo:loop-estimates}
 \mathcal L_T(h)\leq a_phQ_T,\qquad
 \sup_{0\leq h\leq17h_T}\|\mathcal H_T^{(18)}(h)\|\leq Q_T.
\end{equation}
Moreover,
\begin{equation}\label{holo:second-coefficient}
 \mathcal H_T(0)=\Id,\qquad
 \frac12\mathcal H_T''(0)=\pm a_p^2R^{g'}(E_1,E_2)(p_T)
\end{equation}
as endomorphism matrices in the based frame.
\end{lemma}

\begin{figure}[!htbp]
\centering
\begin{tikzpicture}[
  font=\small,
  boundary/.style={draw=black!85, line width=.8pt},
  orient/.style={-{Stealth[length=4pt]}, draw=black!85, line width=.8pt},
  mesh/.style={draw=black!25, line width=.35pt}
]
\node at (8.02,4.25)
  {$\displaystyle\Phi(\sigma,\vartheta)
    =\operatorname{Fl}^{1}_{\sigma f_1+\vartheta f_2}(p_T),
    \qquad f_i=a_pE_i$};
\node[font=\small\bfseries] at (2.28,3.55) {Parameter square};
\node[font=\small\bfseries] at (12.52,3.55) {Image surface (schematic)};

\coordinate (o) at (.80,.30);
\coordinate (a) at (3.50,.30);
\coordinate (b) at (3.50,3.00);
\coordinate (c) at (.80,3.00);
\fill[black!3] (o) rectangle (b);
\draw[mesh] (1.70,.30) -- (1.70,3.00);
\draw[mesh] (2.60,.30) -- (2.60,3.00);
\draw[mesh] (.80,1.20) -- (3.50,1.20);
\draw[mesh] (.80,2.10) -- (3.50,2.10);
\draw[boundary] (o) -- (a) -- (b) -- (c) -- cycle;
\draw[orient] (o) -- (2.32,.30);
\draw[orient] (a) -- (3.50,1.82);
\draw[orient] (b) -- (1.98,3.00);
\draw[orient] (c) -- (.80,1.48);
\fill (o) circle (1.6pt);
\node[anchor=north east] at (o) {$(0,0)$};
\node[anchor=north west] at (a) {$(h,0)$};
\node[anchor=south west] at (b) {$(h,h)$};
\node[anchor=south east] at (c) {$(0,h)$};
\node[fill=white, inner sep=3pt] at (2.15,1.65) {$[0,h]^2$};

\draw[-{Stealth[length=6pt]}, thick] (4.63,1.77) -- (8.65,1.77)
  node[midway, above=5pt] {$\Phi$}
  node[midway, below=6pt, align=center, text width=3.85cm]
    {time-one flow of\\the field\\
     $\sigma f_1+\vartheta f_2$};

\coordinate (p) at (9.72,.37);
\coordinate (q) at (14.40,.65);
\coordinate (r) at (14.67,2.85);
\coordinate (s) at (10.17,3.00);
\path[fill=black!3]
  (p) .. controls (11.15,.80) and (12.85,.10) .. (q)
      .. controls (14.13,1.31) and (15.18,2.03) .. (r)
      .. controls (13.28,3.27) and (11.75,2.53) .. (s)
      .. controls (10.62,2.15) and (9.37,1.32) .. cycle;
\draw[mesh] (11.23,.56)
  .. controls (11.02,1.25) and (11.88,2.24) .. (11.64,2.84);
\draw[mesh] (12.84,.40)
  .. controls (12.58,1.30) and (13.51,2.24) .. (13.17,3.01);
\draw[mesh] (9.72,1.20)
  .. controls (11.20,1.72) and (13.01,.84) .. (14.45,1.38);
\draw[mesh] (10.09,2.12)
  .. controls (11.68,2.52) and (13.09,1.71) .. (14.83,2.11);
\draw[orient] (p) .. controls (11.15,.80) and (12.85,.10) .. (q);
\draw[orient] (q) .. controls (14.13,1.31) and (15.18,2.03) .. (r);
\draw[orient] (r) .. controls (13.28,3.27) and (11.75,2.53) .. (s);
\draw[orient] (s) .. controls (10.62,2.15) and (9.37,1.32) .. (p);
\fill (p) circle (1.8pt);
\node[anchor=north, inner sep=5pt] at (10.26,.37) {$p_T=\Phi(0,0)$};
\node[fill=white, inner sep=3pt] at (12.28,1.72) {$\Phi([0,h]^2)$};
\node[anchor=north] at (12.51,-.24) {$\ell_h=\Phi(\partial[0,h]^2)$};

\node[align=center] at (8.02,-1.26)
  {$\displaystyle\frac12\mathcal H_T''(0)
     =\pm a_p^2R^{g'}(E_1,E_2)(p_T),
   \qquad \mathcal L_T(h)\le a_phQ_T,
   \quad Q_T=e^{o(T)}$.};
\end{tikzpicture}
\caption{The loop surface of Equation~\eqref{holo:surface}, written on
$[0,h]^2$ rather than the unit square. Each image point is the time-one
flow of the single vector field $\sigma f_1+\vartheta f_2$ from $p_T$;
the oriented boundary gives the based loop. The drawing is schematic,
with $a_p$ constant on the observation box. Lemma~\ref{holo:loops}
gives the quadratic holonomy coefficient (up to the stated sign
convention) and the developed length in the transported test metric.}
\label{fig:parameter-flow-loop}
\end{figure}


\begin{proof}
Finite products and sums of subexponential bounds remain
subexponential; $Q_T$ may be increased finitely many times below.
Use $r\in[0,1]$ as the flow variable.  With
$c^\alpha=h(s\delta_1^\alpha+t\delta_2^\alpha)$, the flow solves
\begin{equation}\label{holo:flow-ode}
 \partial_r\Psi=c^\alpha f_\alpha(\Psi),\qquad \Psi(0)=p_T.
\end{equation}
Its coordinate speed and its position-Lipschitz coefficient are
bounded by $hQ_T$.  Since $17h_TQ_T\to0$, the speed bound keeps
the flow in the fixed strict box for the entire interval $[0,1]$;
local existence and continuation give the stated time-one flow.

Coordinate parameter derivatives are also controlled.  A derivative
of total order $n\geq1$ in $(h,s,t)$ satisfies a linear equation
whose coefficient on its order-$n$ unknown is
$c^\alpha Df_\alpha(\Psi)$, of size $hQ_T$.  Its other terms are
finite products of derivatives of $f$ through order $n$, lower
derivatives of $\Psi$, and derivatives of the polynomial $c$.
The initial derivatives vanish because $p_T$ is parameter
independent.  Induction through $n=19$, using
$\exp(O(hQ_T))\leq2$ for large $T$, bounds each such derivative
by $e^{o(T)}$.  This uses at most nineteen coordinate derivatives
of $f$.

The coordinate bounds alone do not give the desired physical
length.  To keep the factor $a_p$, consider the $s$ variation in
the $f$ frame.  Define $J_s(r)$ by the linear system
\begin{equation}\label{holo:variation}
 \frac{dJ_s^\gamma}{dr}
   =h\delta_1^\gamma
       +c^\alpha J_s^\beta
                    B^\gamma{}_{\beta\alpha}(\Psi(r)),
       \qquad J_s(0)=0.
\end{equation}
Then $J_s^\beta f_\beta(\Psi)$ equals $\partial_s\Psi$:
differentiating this vector using Equation~\eqref{holo:variation}
and the definition of the bracket cancels the frame derivative
terms and gives exactly the coordinate variation equation of
Equation~\eqref{holo:flow-ode}, with the same zero initial value.
Uniqueness proves the equality.  This construction never estimates
an inverse coordinate frame.  The equation for $J_t$ has direct
source $h\delta_2^\gamma$ and zero initial value.  In particular
these are inhomogeneous equations for the parameter-dependent
flow; their homogeneous matrix has size $hQ_T$.  Integration gives
\begin{equation}\label{holo:variation-bound}
 |J_s|+|J_t|\leq hQ_T.
\end{equation}

Differentiate Equation~\eqref{holo:variation} directly in $h$,
including at $h=0$.  At derivative order $n$, the coefficient on
$\partial_h^nJ$ is still $c^\alpha B_{\beta\alpha}(\Psi)$.
Every other term contains a lower $h$ derivative of $J$, already
bounded coordinate derivatives of $\Psi$, and derivatives of $B$
of order at most $n$.  The direct source has first derivative
$\delta_i$ and all higher derivatives zero.  Thus induction gives
\[
 \sup_{r,s,t,h}|\partial_h^nJ_s|
     +|\partial_h^nJ_t|\leq e^{o(T)},\qquad 1\leq n\leq18.
\]
No division by $h$ is used.  The twenty-derivative bounds for
$f,B,G$ allow the nineteen coordinate variation derivatives and
the eighteen transport derivatives with a spare derivative.

Parametrize each oriented boundary edge by $z\in[0,1]$.
Its tangent has the form $J^\alpha f_\alpha$, where $J$ is one
of $\pm J_s$, $\pm J_t$ evaluated at $r=1$ and at that edge.
Parallel transport in the moving $E$ frame solves
\begin{equation}\label{holo:transport}
 \frac{dP}{dz}=-\mathcal G(z,h)P,\qquad
 \mathcal G^\mu{}_{\beta}(z,h)
          =G^\mu{}_{\alpha\beta}(\Psi_h)J^\alpha.
\end{equation}
The undifferentiated coefficient is $O(hQ_T)$ by
Equation~\eqref{holo:variation-bound}; its first eighteen
$h$ derivatives have size $e^{o(T)}$.  Fundamental matrices and
their inverses are bounded by $\exp(O(hQ_T))$ on each edge.
Differentiating Equation~\eqref{holo:transport} $n$ times yields
the same homogeneous coefficient on $\partial_h^nP$ and source
\[
 -\sum_{j=1}^n\binom nj
          (\partial_h^j\mathcal G)(\partial_h^{n-j}P).
\]
Induction through $n=18$, followed by multiplication of the four
edge matrices, gives the derivative bound in
Equation~\eqref{holo:loop-estimates}.

At any point on the boundary the tangent has $E$-components
$a_pJ$.  If $P$ denotes accumulated parallel transport from the
base, its norm in the loop-parallel positive metric is its base
norm after applying $P^{-1}$, hence at most
$C a_p|P^{-1}J|\leq a_phQ_T$ by the uniform base comparison and
Equation~\eqref{holo:variation-bound}.  Summing the four unit edge
intervals proves the length estimate in the actual test metric.

Finally consider the two-parameter surface
$\Phi(\sigma,\vartheta)=
\operatorname{Fl}^1_{\sigma f_1+\vartheta f_2}(p_T)$.
Our loops are its boundaries of $[0,h]^2$.  For each fixed $T$,
pull back the connection and choose a smooth frame equal to
$E(p_T)$ at the origin whose connection one-form vanishes there.
Such a frame is obtained by prescribing its first derivative at
the origin; no uniform estimate on that auxiliary choice is used.
Its connection one-form is $O(|\sigma|+|\vartheta|)$.  Therefore
ordered products of two edge integrals have order at least $h^4$,
while Stokes' formula gives the order-$h^2$ term as minus the
curvature at the origin.  The differential of $\Phi$ there sends
the coordinate vectors to $f_1(p_T),f_2(p_T)$.  This proves
Equation~\eqref{holo:second-coefficient}, with the sign accounting
for orientation and curvature convention.  The uniform derivative
bound was proved separately in the original frame.
\end{proof}

\subsection{An exact eighteen-node formula}

\begin{lemma}[Extraction of the quadratic coefficient]
\label{holo:interpolation}
Let $F:[0,17h]\to V$ be $C^{18}$, where $V$ is a
finite-dimensional normed vector space and $h>0$.  Put
\[
 S_1=\sum_{r=1}^{17}\frac1r,\qquad
 S_2=\sum_{r=1}^{17}\frac1{r^2},\qquad
 c_0=\frac{S_1^2-S_2}{2},\qquad
 c_j=\frac{(-1)^j}{j}\binom{17}{j}
                   \left(S_1-\frac1j\right)\quad(1\leq j\leq17).
\]
Then the following identity holds with Bochner integrals in $V$:
\begin{equation}\label{holo:exact-interpolation}
 \frac{F''(0)}2
  =h^{-2}\sum_{j=0}^{17}c_jF(jh)
    -\frac{h^{-2}}{17!}\sum_{j=0}^{17}c_j
       \int_0^{jh}(jh-t)^{17}F^{(18)}(t)\,dt.
\end{equation}
In particular, if $F(0)=0$, $\|F(jh)\|\leq b jh$, and
$\sup\|F^{(18)}\|\leq M$, then
\begin{equation}\label{holo:interpolation-bound}
 \left\|\frac{F''(0)}2\right\|
       \leq K_1\frac b h+K_2 M h^{16},\qquad
 K_1=\sum_{j=1}^{17}|c_j|j,\quad
 K_2=\frac1{18!}\sum_{j=1}^{17}|c_j|j^{18}.
\end{equation}
Both constants are finite and independent of $F,h,T$.
\end{lemma}

\begin{proof}
Let $L_j(x)=\prod_{0\leq r\leq17,\ r\ne j}(x-r)/(j-r)$.
The coefficient of $x^2$ in $L_0$ is
$\sum_{1\leq r<s\leq17}(rs)^{-1}=c_0$.  For $j\geq1$,
$L_j$ has a factor $x$,
$L_j'(0)=(-1)^{j-1}\binom{17}{j}/j$, and the logarithmic
derivative at zero of its remaining product is $-S_1+1/j$.
The coefficient of $x^2$ is therefore $c_j$ as displayed.
Lagrange interpolation gives the exact moment identities
\[
 \sum_{j=0}^{17}c_jj^k=
     \begin{cases}1,&k=2,\\0,&0\leq k\leq17, k\ne2.
     \end{cases}
\]
Apply these identities to the Taylor polynomial
$\sum_{k=0}^{17}F^{(k)}(0)t^k/k!$, evaluated at $t=jh$.
Taylor's integral remainder gives
Equation~\eqref{holo:exact-interpolation}.  Bounding each remainder
by $M(jh)^{18}/18!$ gives
Equation~\eqref{holo:interpolation-bound}.
\end{proof}

\begin{proposition}[Finite small-loop violation]
\label{holo:violation}
Fix a passing test $\mathcal A(O,B)$, a datum
$d\in\mathcal A(O,B)$, and a branch seed $w\in O$ as in
Corollary~\ref{test:branch-seed}.  For every finite $m\geq10$ the
signed experiments can be chosen with
$p_m(d_T^\pm-d)\to0$ such that, for every sufficiently large $T$,
at least one of $d_T^+$ and $d_T^-$ does not belong to
$\mathcal A(O,B)$.
\end{proposition}

\begin{proof}
Choose the sign in Proposition~\ref{holo:signal}.  Its adaptive
seed $w_T$ is in $O$ for large $T$ and reaches $p_T$ in the full
MGHD at a finite positive proper time, by
Proposition~\ref{attach:observer}.  The loops in
Lemma~\ref{holo:loops} lie entirely in $I^+(\Sigma)$ of that
development by Proposition~\ref{attach:development}.
Their lengths tend to zero uniformly for $0\leq h\leq17h_T$,
so they satisfy the strict cutoff $\mathcal L_T(h)<B^{-1}$.

Suppose the selected datum still belongs to $\mathcal A(O,B)$.
Its loop inequality in Definition~\ref{test:definition}, the
uniform frame comparison, and Equation~\eqref{holo:loop-estimates}
then give
\[
 \|\mathcal H_T(jh_T)-\Id\|
          \leq a_p jh_T e^{o(T)},\qquad 0\leq j\leq17.
\]
The fixed constants, including $B$, have been included in
$e^{o(T)}$.  Apply Lemma~\ref{holo:interpolation} with
$F=\mathcal H_T-\Id$, $b=a_pe^{o(T)}$, and $M=e^{o(T)}$.
Equation~\eqref{holo:second-coefficient} gives
\begin{align}
 \|a_p^2R^{g'}(E_1,E_2)(p_T)\|
 &\leq e^{o(T)}(a_p/h_T+h_T^{16})                         \label{holo:curvature-upper}\\
 &\leq e^{o(T)}
        (e^{-15\kappa T/16}+e^{-\kappa T})=o(A_T).\notag
\end{align}
Indeed $A_T=e^{-14\kappa T/16}$, so the two ratios have exponents
$-\kappa T/16+o(T)$ and $-\kappa T/8+o(T)$, respectively.
Both are strict, independent of the fixed interpolation constants.

Because the Lorentz Gram matrix of $E$ is the constant matrix
$\eta$, the sectional component
$a_p^2\Riem(g')(E_1,E_2,E_1,E_2)$ is a fixed matrix contraction
of the endomorphism in Equation~\eqref{holo:curvature-upper}.
It is therefore $o(A_T)$, contradicting $|C_T^\pm|>A_T$.
The chosen sign fails the test.  The convergence in the prescribed
finite seminorm and membership in the admissible data neighborhood
are those of the signed construction; hence both data can be
placed in any prescribed relative neighborhood of $d$.
\end{proof}
